% =====================================================================
% 系统开发工具基础 · 实验报告（第 13--16 题）
%
% 编译方式（在 WSL 中）：
%   latexmk -xelatex -halt-on-error 实验报告.tex
% =====================================================================
\documentclass[UTF8,fontset=ubuntu]{ctexart}

\usepackage{amsmath}
\usepackage{booktabs}
\usepackage{geometry}
\usepackage{listings}
\usepackage[hidelinks]{hyperref}

\geometry{a4paper, margin=2.5cm}

\lstset{
  basicstyle=\ttfamily\small,
  breaklines=true,
  frame=single,
  framerule=0.4pt,
  tabsize=4,
  showstringspaces=false,
  columns=fullflexible,
  keepspaces=true,
}

\title{\textbf{系统开发工具基础 · 实验报告（第 13--16 题）}}
\author{姓名：李云龙 \quad 学号：25030021043}
\date{\today}

\begin{document}

\maketitle

\begin{center}
\begin{tabular}{@{}p{3cm}p{11.5cm}@{}}
\toprule
姓名 & 李云龙 \\
学号 & 25030021043 \\
实验日期 & \today \\
实验环境 & Windows 11 + WSL2（Ubuntu 26.04 LTS）；Python 3.14 + ruff 0.16 + pytest 9.1 + GNU Make 4.4.1 + jq 1.8.1 \\
\bottomrule
\end{tabular}
\end{center}

\begin{quote}
\textbf{说明：}本报告中各题“运行结果”均为在 WSL 中实际运行的输出；第 16 题的“失败门禁”输出与 wheel 的 SHA-256 均如实记录。
\end{quote}

\section{实验目的}
\begin{enumerate}
\item 建立一个可执行的本地质量门禁（格式检查、静态检查、测试三合一），并理解“不全局忽略规则”的意义。
\item 理解 Make 的增量构建判据（目标与依赖的时间戳），写出依赖精确的 Makefile。
\item 掌握用 curl + jq 把本地 API 的 JSON 数据转换为可读 Markdown 报告。
\item 在陌生仓库中完成一次综合检查：用门禁暴露缺陷、定位修复、产出可校验的构建产物。
\end{enumerate}

\section{实验十三：建立可执行的本地质量门禁（q13）}

\subsection{最小配置（pyproject.toml）}
\begin{lstlisting}
[build-system]
requires = ["setuptools>=68"]
build-backend = "setuptools.build_meta"

[project]
name = "greetlab-25030021043"
version = "0.1.0"

[project.scripts]
sdt-greet = "greetlab.cli:main"

[tool.ruff]
line-length = 88
src = ["src", "tests"]

[tool.pytest.ini_options]
testpaths = ["tests"]
\end{lstlisting}

只保留最小配置，\textbf{没有使用任何全局忽略规则}（如 \lstinline|ignore = [...]|）。

\subsection{测试（正常姓名 + 空白姓名，均有有效断言）}
\begin{lstlisting}[language=Python]
import sys

import pytest

from greetlab.cli import main


def test_greet_normal_name(monkeypatch, capsys):
    monkeypatch.setattr(sys, "argv", ["sdt-greet", "--name", "25030021043"])
    main()
    captured = capsys.readouterr()
    assert captured.out == "Hello, 25030021043!\n"


@pytest.mark.parametrize("blank", [" ", "   ", "\t"])
def test_blank_name_exits_with_code_2(monkeypatch, capsys, blank):
    monkeypatch.setattr(sys, "argv", ["sdt-greet", "--name", blank])
    with pytest.raises(SystemExit) as excinfo:
        main()
    assert excinfo.value.code == 2
    assert capsys.readouterr().out == ""
\end{lstlisting}

\subsection{质量门禁脚本（check.sh）}
\begin{lstlisting}[language=bash]
#!/usr/bin/env bash
# 本地质量门禁：格式检查 -> 静态检查 -> 测试，任一环节失败即返回非 0
set -euo pipefail
cd "$(dirname "$0")"

echo "== 1/3 ruff format --check =="
ruff format --check .

echo "== 2/3 ruff check =="
ruff check .

echo "== 3/3 pytest =="
python3 -m pytest

echo "全部通过（exit 0）"
\end{lstlisting}

\subsection{运行结果}
\begin{lstlisting}
$ ruff format .
4 files left unchanged

$ ./check.sh
== 1/3 ruff format --check ==
4 files already formatted
== 2/3 ruff check ==
All checks passed!
== 3/3 pytest ==
tests/test_cli.py ....                                                   [100%]
4 passed in 0.26s
全部通过（exit 0）
\end{lstlisting}

\subsection{要点分析}
\begin{itemize}
\item 门禁把“格式 + 静态检查 + 测试”固化为\textbf{一条可执行命令}，本地开发与 CI 使用同一标准，避免“我这儿能过”；
\item \lstinline|set -euo pipefail| 保证任一环节失败立即中止并返回非 0，门禁才有“门”的作用；
\item \textbf{不全局忽略规则}是关键要求：一旦写 \lstinline|ignore = [...]|，等于把规则悄悄关掉，门禁形同虚设；
\item 两个测试各自带有效断言：正常路径断言\textbf{输出内容}，异常路径断言\textbf{退出码}并额外确认“没有输出问候语”。
\end{itemize}

\section{实验十四：让 Make 只重建真正受影响的产物（q14）}

\subsection{源文件（题目给定）}
\textbf{data.csv}：
\begin{lstlisting}
name,value
alpha,2
beta,3
gamma,5
\end{lstlisting}

\textbf{stats.py}：
\begin{lstlisting}[language=Python]
import csv

with open("data.csv") as f:
    total = sum(int(r["value"]) for r in csv.DictReader(f))
open("stats.txt", "w").write(str(total))
\end{lstlisting}

\textbf{build\_report.py}：
\begin{lstlisting}[language=Python]
text = open("report.md").read()
total = open("stats.txt").read()
open("report.txt", "w").write(f"{text}\nTotal: {total}\n")
\end{lstlisting}

\textbf{report.md}：
\begin{lstlisting}
# Data Report
\end{lstlisting}

\subsection{Makefile}
\begin{lstlisting}
.PHONY: all clean

all: stats.txt report.txt

stats.txt: data.csv stats.py
	python3 stats.py

report.txt: report.md stats.txt build_report.py
	python3 build_report.py

clean:
	rm -f stats.txt report.txt
\end{lstlisting}

\subsection{验证结果}
\begin{center}
\begin{tabular}{@{}p{5.5cm}p{9cm}@{}}
\toprule
操作 & 实际现象 \\
\midrule
首次 \lstinline|make| & 依次执行 stats.py、build\_report.py；stats.txt 内容为 10，report.txt 为 Data Report + Total: 10 \\
无改动再 \lstinline|make| & \lstinline|make: Nothing to be done for 'all'.|（不执行任何配方） \\
\lstinline|touch data.csv| 后 \lstinline|make| & 两个配方都执行（data.csv 变新 -> stats.txt 重建 -> 进而 report.txt 重建） \\
只 \lstinline|touch report.md| 后 \lstinline|make| & \textbf{只执行 build\_report.py}，不牵连上游 \\
\lstinline|make clean| & 只删除 stats.txt、report.txt；源文件与 Makefile 均保留 \\
\bottomrule
\end{tabular}
\end{center}

\subsection{要点分析}
\begin{itemize}
\item Make 的重建判据是“\textbf{目标文件与依赖文件的时间戳}”：任一依赖比目标新，则执行配方；
\item 因此依赖必须\textbf{列全且精确}：stats.txt 依赖 data.csv 与 stats.py，report.txt 依赖 report.md、stats.txt 与 build\_report.py。漏列会漏建（改脚本却不重建），多列会白建（无关改动触发重建）；
\item \lstinline|all|、\lstinline|clean| 不是磁盘上的真实文件，必须声明 \lstinline|.PHONY|，否则一旦目录里出现同名文件，Make 会用时间戳判定“已是最新”而跳过配方。
\end{itemize}

\section{实验十五：把本地 API 数据转换为可读报告（q15）}

\subsection{数据与脚本}
packages.json（题目给定，5 条记录）保存于 q15 目录，用 \lstinline|python3 -m http.server 8000| 提供本地 HTTP 访问。\textbf{api\_report.sh}：
\begin{lstlisting}[language=bash]
#!/usr/bin/env bash
# 从本地 HTTP 服务拉取 packages.json，用 jq 筛选排序，生成 Markdown 报告 summary.md
set -euo pipefail
cd "$(dirname "$0")"

url="${1:-http://127.0.0.1:8000/packages.json}"
tmp_json="$(mktemp)"
trap 'rm -f "$tmp_json"' EXIT

if ! curl -fsS "$url" -o "$tmp_json"; then
    echo "错误：无法从 $url 获取数据（请确认已在 q15 目录启动 python3 -m http.server 8000）" >&2
    exit 1
fi

{
    echo "# Package Summary"
    echo
    echo "数据来源：${url}；筛选条件：status=active 且 downloads >= 100；排序：downloads 降序，name 升序。"
    echo
    jq -r '["| name | version | downloads |", "| --- | --- | --- |"]
           + ([.[] | select(.status == "active" and .downloads >= 100)]
              | sort_by(-.downloads, .name)
              | map("| \(.name) | \(.version) | \(.downloads) |"))
           | .[]' "$tmp_json"
} > summary.md

cat summary.md
\end{lstlisting}

\subsection{运行结果（summary.md）}
\begin{lstlisting}
# Package Summary

数据来源：http://127.0.0.1:8000/packages.json；筛选条件：status=active 且 downloads >= 100；排序：downloads 降序，name 升序。

| name | version | downloads |
| --- | --- | --- |
| delta | 0.9.0 | 450 |
| gamma | 1.5.1 | 450 |
| alpha | 1.2.0 | 120 |
\end{lstlisting}

服务端日志确认请求真实发生：
\begin{lstlisting}
127.0.0.1 - - [14/Sep/2026 14:01:22] "GET /packages.json HTTP/1.1" 200 -
\end{lstlisting}

\subsection{要点分析}
\begin{itemize}
\item 筛选结果正确：beta（status=inactive）与 epsilon（downloads=80）被排除；
\item 排序正确：delta 与 gamma 同为 450，按 name 升序时 delta 在前；alpha 120 最后；
\item \lstinline|curl -fsS| 的组合意义：\lstinline|-f| 让 HTTP 4xx/5xx 也返回非零退出码（配合 \lstinline|set -e| 及时中止）、\lstinline|-s| 静默进度条、\lstinline|-S| 保留错误信息；
\item jq 用一条表达式完成“筛选 + 多键排序”（\lstinline|select| + \lstinline|sort_by(-.downloads, .name)|），比用 shell 循环处理 JSON 更可靠。
\end{itemize}

\section{实验十六：修复并交付一个陌生的小型工具仓库（q16）}

\subsection{Makefile}
\begin{lstlisting}
.PHONY: check build clean

check:
	./check.sh

build:
	python3 -m build

clean:
	rm -rf build dist src/*.egg-info
\end{lstlisting}

\subsection{综合检查过程}
\textbf{第 1 步：基线检查}
\begin{lstlisting}
$ make check
4 passed in 0.23s
全部通过（exit 0）
\end{lstlisting}

\textbf{第 2 步：把问候语临时改成字面量}（模拟“坏实现”）
\begin{lstlisting}[language=Python]
    print("Hello, name!")      # 无论传入什么姓名，都输出固定的字面量
\end{lstlisting}

\textbf{第 3 步：门禁成功拦住}
\begin{lstlisting}
$ make check
>       assert captured.out == "Hello, 25030021043!\n"
E       AssertionError: assert 'Hello, name!\n' == 'Hello, 25030021043!\n'
FAILED tests/test_cli.py::test_greet_normal_name
1 failed, 3 passed in 0.21s
make: *** [Makefile:4: check] Error 1
\end{lstlisting}

\textbf{第 4 步：定位并修复}（恢复 f-string，使用真实参数）
\begin{lstlisting}[language=Python]
    print(f"Hello, {a.name}!")
\end{lstlisting}

\textbf{第 5 步：复验与构建}
\begin{lstlisting}
$ make check
4 passed in 0.17s
全部通过（exit 0）

$ make build
Successfully built greetlab_25030021043-0.1.0.tar.gz and greetlab_25030021043-0.1.0-py3-none-any.whl
\end{lstlisting}

\textbf{第 6 步：计算 wheel 的 SHA-256}
\begin{lstlisting}
09f86db07ed755d4dbb0e73f8746527a843a7950c22504969812ed05335ff8e3  dist/greetlab_25030021043-0.1.0-py3-none-any.whl
\end{lstlisting}

\textbf{第 7 步：提交}——最终改动以内容聚焦的 Git 提交入库（提交信息中记录了 wheel 的 SHA-256 前 8 位）。

\subsection{要点分析}
\begin{itemize}
\item 一个“陌生仓库”能不能放心修改，取决于\textbf{是否存在可执行的质量门禁}：测试失败直接导致 \lstinline|make check| 失败（退出码非 0），说明门禁真的在起作用；
\item 修复遵循最小改动原则（只恢复被改坏的一行），修改后用\textbf{同一条命令}复验，避免引入新的偏差；
\item “可交付”的标志是产物可校验：生成 wheel 并给出 SHA-256，接收方可用同一命令重建并比对哈希。
\end{itemize}

\section{遇到的问题与解决}
\begin{center}
\begin{tabular}{@{}p{4.6cm}p{5.2cm}p{5.4cm}@{}}
\toprule
问题 & 原因 & 解决 \\
\midrule
make / jq 命令不存在 & Ubuntu 基础镜像未预装 & apt install make jq \\
测试导入不到 src 布局的包 & src 布局下包不在 sys.path & 增加 conftest.py 把 src 加入 sys.path \\
启动本地服务后又“消失” & pkill -f http.server 的正则匹配到执行该命令的 shell 自身，把自己杀掉 & 改为脚本内用 \$! 记录 PID + trap 收尾（demo.sh） \\
坏实现是否会被发现 & 需要可执行门禁验证 & 故意改成字面量后 make check 立即失败，证明门禁有效 \\
\bottomrule
\end{tabular}
\end{center}

\section{实验总结}
通过本次实验，我建立了由 ruff format、ruff check 与 pytest 组成的本地质量门禁，并体会到“门禁必须可执行、且不能靠全局忽略规则绕过”；理解了 Make 依据时间戳做增量构建的原理，写出依赖精确、all/clean 均为 .PHONY 的 Makefile，并验证了“只改下游输入则只重建下游产物”；掌握了用 curl -fsS 与 jq 把本地 HTTP 服务上的 JSON 转换为带筛选、多键排序的 Markdown 报告；最后在陌生仓库中完成综合检查——先用门禁暴露缺陷、最小修复、再复验并产出带 SHA-256 的 wheel，完整体验了可交付的工程闭环。

\section*{附录：文件清单}
\begin{lstlisting}
w4/README.md                  第 13--16 题说明（步骤、原理、实测数据）
w4/q13/check.sh               第 13 题质量门禁脚本
w4/q13/pyproject.toml         第 13 题 ruff / pytest 最小配置
w4/q13/tests/test_cli.py      第 13 题测试（正常姓名 + 空白姓名）
w4/q13/src/greetlab/          第 13 题包源码
w4/q14/Makefile               第 14 题构建规则（依赖精确 + .PHONY）
w4/q14/{data.csv,stats.py,build_report.py,report.md}   第 14 题源文件
w4/q14/{stats.txt,report.txt} 第 14 题生成的产物
w4/q15/packages.json          第 15 题数据
w4/q15/api_report.sh          第 15 题报告生成脚本
w4/q15/demo.sh                第 15 题一键演示（起服务 + 生成 + 收尾）
w4/q15/summary.md             第 15 题生成的 Markdown 报告
w4/q16/Makefile               第 16 题 check / build / clean
w4/q16/check.sh               第 16 题质量门禁（同 q13）
w4/q16/src/greetlab/cli.py    第 16 题修复后的实现
w4/q16/dist/                  第 16 题构建产物（wheel，SHA-256 见 5.2）
实验报告.md / 实验报告.tex / 实验报告.pdf    本报告
\end{lstlisting}

\end{document}
